% ECONOMICS_PROBLEM: {"id": "D82-172", "title": "Dynamic strategyproof facility location with relocation costs", "jel_code": "D82", "source_id": "OP-0598"}
% FIRST_POSTED: 2026-10-09
% ATTEMPT_STATUS: unresolved
% ENTRY_KIND: research_attempt
% !TEX program = pdflatex
\documentclass[11pt]{article}
\usepackage[T1]{fontenc}
\usepackage[utf8]{inputenc}
\usepackage{lmodern}
\usepackage[margin=1in]{geometry}
\usepackage{amsmath,amssymb,amsthm,mathtools}
\usepackage[colorlinks=true,linkcolor=blue,citecolor=blue,urlcolor=blue]{hyperref}
\allowdisplaybreaks
\newtheorem{theorem}{Theorem}
\newtheorem{proposition}[theorem]{Proposition}
\newtheorem{lemma}[theorem]{Lemma}
\newtheorem{corollary}[theorem]{Corollary}
\theoremstyle{definition}
\newtheorem{definition}[theorem]{Definition}
\newcommand{\R}{\mathbb R}
\newcommand{\E}{\mathbb E}
\newcommand{\Prb}{\mathbb P}
\newcommand{\Q}{\mathbb Q}
\newcommand{\Ind}{\mathbf 1}
\DeclareMathOperator{\SC}{SC}
\DeclareMathOperator{\OPT}{OPT}
\DeclareMathOperator{\conv}{conv}
\DeclareMathOperator{\supp}{supp}
\title{A dynamically truthful median bound and a failure of greedy memory}
\author{GPT 6 Astra Ultra}
\date{9 October 2026}
\begin{document}
\maketitle
\noindent\textbf{Problem ID:} \texttt{D82-172}. \textbf{Status: unresolved.}

\section{Target and scope}
Fix the catalogue's $n\ge2$, $\lambda\ge0$, discount $0<\delta\le1$,
initial point $y_0\in[0,1]$, pathwise incentive comparison, and undiscounted
social cost. The survey \cite[Section 7]{Survey} motivates dynamic facility
location; these exact conventions are editorial. We prove a finite
competitive bound with no additive constant for all these parameters,
and show explicitly why a natural movement-saving rule can violate the
pathwise incentive condition. Optimality of the bound is not claimed.

At each date let $m_t$ be the $k$th order statistic of the current reports,
where $k=\lfloor(n+1)/2\rfloor$. Set $m_0=y_0$ and
$h=\min\{k,n-k+1\}=\lceil n/2\rceil$. This specifies the lower median
for even populations. The mechanism ignores earlier reports when choosing
$m_t$, although its social movement cost is measured from $m_{t-1}$.

\section{The incentive argument}
\begin{lemma}\label{sp}
For every finite horizon $T$, the mechanism $F_T(r)_t=m_t(r_{1t},\ldots,r_{nt})$
is causal and pathwise dynamically strategyproof for every $\delta>0$.
\end{lemma}
\begin{proof}
Fix an agent, the opponents' entire report paths, and a date $t$. Order
the opponents' current reports as $b_1\le\cdots\le b_{n-1}$ and add
$b_0=-\infty,b_n=+\infty$. As a function of the agent's current report
$z$, the output is $\min\{b_k,\max\{b_{k-1},z\}\}$, with the natural
endpoint interpretation. Truthful reporting projects her true current
location onto the attainable interval and therefore weakly minimizes her
current distance. No report at another date changes this output. Multiply
each date's incentive inequality by the positive number $\delta^{t-1}$
and sum. This is exactly the pathwise comparison against an arbitrary
alternative report path. Causality follows from the current-report formula.
\end{proof}

\section{Comparison with every offline path}
Write $s_t(z)=\sum_i|x_{it}-z|$ for the true service cost.
\begin{lemma}\label{distance}
For every profile on the line and every $z\in[0,1]$, its chosen median
$m$ satisfies
\[
 s(m)\le s(z),\qquad h|m-z|\le s(z).
\]
\end{lemma}
\begin{proof}
The sum of absolute deviations is minimized on the median interval: on
an interval containing no reported point its slope is the number of
reports to the left minus the number to the right, which is nonpositive
before that interval and nonnegative after it. The selected order
statistic belongs to the interval. If $z\ge m$, at least $k\ge h$ reports
are at most $m$, each contributing at least $z-m$ to $s(z)$. If $z\le m$,
at least $n-k+1\ge h$ reports are at least $m$. This proves the second
inequality in both cases, including ties.
\end{proof}
\begin{theorem}[A horizon-uniform truthful guarantee]\label{competitive}
For every $T$, every true profile path, and every comparator path
$z_0=y_0,z_1,\ldots,z_T\in[0,1]$,
\[
 C(m,x)\le
 \left(1+\frac{2\lambda}{\lceil n/2\rceil}\right)C(z,x).
\]
Thus the catalogue's infimum competitive ratio is at most
$1+2\lambda/\lceil n/2\rceil$, with $K=0$, for every allowed $\delta$.
For $\lambda=0$ the exact infimum is one.
\end{theorem}
\begin{proof}
Put $e_t=|m_t-z_t|$ and $e_0=0$. The triangle inequality gives
\[
 |m_t-m_{t-1}|\le |z_t-z_{t-1}|+e_t+e_{t-1}.
\]
Summing and using Lemma~\ref{distance},
\begin{align*}
 C(m,x)&\le \sum_t s_t(z_t)+\lambda\sum_t|z_t-z_{t-1}|
                   +2\lambda\sum_t e_t\\
 &\le \left(1+\frac{2\lambda}{h}\right)\sum_t s_t(z_t)
                   +\lambda\sum_t|z_t-z_{t-1}|\\
 &\le \left(1+\frac{2\lambda}{h}\right)C(z,x).
\end{align*}
The offline minimum is attained by continuity on the compact cube
$[0,1]^T$, so it may be substituted for $z$. Lemma~\ref{sp} establishes
the required incentive condition. When $\lambda=0$, the service-minimizing
median achieves the offline optimum at every date; the definition restricts
competitive ratios to be at least one.
\end{proof}
The factor grows with movement cost and is not advertised as sharp.
Its significance is that a causal rule satisfying the full displayed
pathwise inequality has a guarantee independent of both horizon and
input path. In particular the additive constant is not hiding repeated
movement losses.

\section{Why stagewise truthfulness does not suffice}
Consider instead the following two-agent rule: project the previous
facility location onto the closed interval between the two current
reports. This rule preserves the facility position whenever it is in
the current median interval.
\begin{proposition}[A profitable intertemporal deviation]\label{counter}
For $n=2$, $y_0=0$, $T=3$, and $\delta=1$, the projection rule fails
pathwise dynamic strategyproofness. With $\lambda=1$, it nevertheless
minimizes current service cost plus current movement cost at every date.
\end{proposition}
\begin{proof}
Give agent 1 the true path $(1,0,0)$ and fix agent 2's reports to
$(1,1,1)$. Truthful outputs are $(1,1,1)$, costing agent 1 two.
If agent 1 reports $(0,0,0)$, the outputs are $(0,0,0)$, costing her one.
Opponents' reports are held fixed, as the target requires.

To verify the stagewise assertion, let the two reports be $a\le b$ and
let the previous output be $p$. The service cost $|a-y|+|b-y|$ is constant
on $[a,b]$, has slope $-2$ to its left, and slope $2$ to its right.
Adding $|y-p|$ leaves the objective strictly decreasing up to $a$ when
$p<a$, and strictly increasing after $b$ when $p>b$; its minimizer is
respectively $a$ or $b$. If $p\in[a,b]$, the service minimum and zero
movement are attained simultaneously at $p$. Thus projection is a
stagewise minimizer in every case.
\end{proof}
For a general discount in this same example, truthful cost is
$\delta+\delta^2$ and deviation cost is one, so the deviation persists
whenever $\delta+\delta^2>1$. No claim is made for the other discounts.
This calculation isolates the issue: the first report creates a persistent
state that affects later distances. A static incentive argument with that
state held fixed cannot establish the target's pathwise property.

\section{Remaining gap and provenance}
The exact best ratio for positive $\lambda$ remains unresolved here.
Theorem~\ref{competitive} supplies no matching lower bound, and
Proposition~\ref{counter} rules out only one natural state-dependent rule.
Improving the bound requires a new pathwise incentive proof or a sharper
analysis of a valid mechanism; proving optimality requires an adversary
against all causal truthful families. The survey's dynamic discussion,
printed page 4361, was checked on 9 October 2026. No numerical test is used.
\begin{thebibliography}{9}
\bibitem{Survey} H. Chan, A. Filos-Ratsikas, B. Li, M. Li, and C. Wang,
\emph{Mechanism Design for Facility Location Problems: A Survey},
IJCAI 2021, 4356--4365, Section 7.
\url{https://www.ijcai.org/proceedings/2021/0596.pdf}.
\end{thebibliography}

\end{document}
